\defmodule {maxper}

This program verifies whether a given generator has maximal period or not.
Integers are represented using the {\tt SuperInteger} type.
The modulus $m$ must be prime.
To verify the maximal period conditions, the factorizations of $m-1$ and
$r = (m^k-1)/(m-1)$ are required.  They can be found by the program
or provided by the user in a file.
We warn that factoring $r$ can take a {\em huge} amount of time.

The data for {\tt maxper} must be placed in a file with extension
``{\tt .dat}'', according to the format displayed in Figure~\ref{maxper.dat}.
The data fields have the same meaning as for {\tt latl}.
To run the program, type ``{\tt maxper \em\<file\>}'', 
where {\em\<file\>} is the name of the data file, without extension.
Currently, only {\tt LCG} and {\tt MRG} are allowed for {\em TypeGen}.

%%  DATA FILE  %%%%%%%%%%%%%%%%%%%%%%%%%%%%%%%%%%%%%%%%
\setbox0=\vbox {\hsize=12.0cm \setverbatim\it\baselineskip=13pt

           TypeGen
           $m$
           $k$
           $F1$ $[$\<file1\>$]$
           $F2$ $[$\<file2\>$]$
           $a_1$
           \vdots
           $a_k$

}
\begin{figure}[h] \centerline {\boxit{\box0}}
\caption { Data file format for {\tt maxper}. }  \label {maxper.dat}
\end{figure}

